% ---------------------------------------------------------------------------
% Drop-in subsection for main.tex, to follow \subsection{Longitudinal response
% at every momentum} (sec:szz).  Requires \Jpm, \Jzz, \bq as already defined,
% and the figure paper/figs/fig_fcc32_parity.pdf.
%
% What it establishes, in order:
%   1. why the eight cluster momenta are three stars and nothing disperses;
%   2. that S^zz misses the lowest excitation at X entirely;
%   3. that the miss is an exact spin-flip selection rule, not a small number,
%      demonstrated by the complementary magnetic channel.
% ---------------------------------------------------------------------------

\subsection{What the longitudinal probe can and cannot see}
\label{sec:parity}

The eight FCC-32 momenta are not eight independent points: they are three
high-symmetry stars of the fcc Brillouin zone -- one $\Gamma$, three $X$, four
$L$ -- and both nonzero stars lie on the zone boundary.  The cluster therefore
carries no small-$\bq$ information at all, and no statement about a photon
\emph{dispersion} can be made on it.  What can be made, and is the subject of
this subsection, is a statement about which excitations the longitudinal probe
is able to reach.

To separate the probe from the spectrum we resolve the eigenvalues by momentum
directly, projecting each eigenvector of the masked block onto the cluster
momenta with $P_{\bq}=N_t^{-1}\sum_t \chi_{\bq}(t)^{*}T_t$, where $T_t$ runs over
the lattice translations.  The projector is exact to $4\times10^{-13}$.  Each
level then carries a degeneracy equal to its star size times two.

The result is that $S^{zz}$ does not couple to the lowest excitation at $X$.
At $\Jpm/\Jzz=-0.20$ the lowest $X$ level lies at $\omega/\Jzz=0.025274$ and
carries longitudinal weight $3.2\times10^{-27}$, while the weight the structure
factor does show at $X$ sits on a higher level at $0.037928$.  The same holds
at $\Jpm/\Jzz=+0.046$: the lowest $X$ level at $0.003612$ carries
$4.6\times10^{-27}$, and the visible weight sits at $0.004525$.  At $L$, by
contrast, the probe sits on the lowest level and takes $88$--$90\%$ of that
star's weight.  Resolving the four sublattice channels separately does not
change this: every level the probe reaches has all four channels exactly equal,
and every level it misses is at machine zero in all four.  The miss is a
selection rule, not a form factor.

\begin{figure*}[t]
\centering
\includegraphics[width=\textwidth]{fig_fcc32_parity.pdf}
\caption{Electric and magnetic probes reach disjoint sets of levels.
(a), (b) Every computed level of the masked block, drawn at its own energy in
the column of its momentum star, coloured by which probe reaches it and with
bar width encoding the weight; levels reached by neither probe are shown as
hairlines so that they remain visible.  (c) The electric weight of each level
against its magnetic weight.  Points pin to one axis or the other over more
than twenty-five decades and the diagonal region, where a level reached by both
probes would sit, is empty.  This is Eq.~(\ref{eq:parity}) at work rather than
a numerical coincidence.}
\label{fig:parity}
\end{figure*}

The rule is spin-flip parity.  The Hamiltonian is invariant under the global
flip $\mathcal{P}: S^z_i\mapsto-S^z_i$, $S^\pm_i\mapsto S^\mp_i$, which
preserves the $\Jzz$ term trivially and maps the $\Jpm$ term onto its Hermitian
conjugate.  The longitudinal probe is odd under it, while the hexagon ring
exchange -- the lattice magnetic flux
$\Phi_h=S^+S^-S^+S^-S^+S^-+\mathrm{h.c.}$ -- is even:
\begin{equation}
 \mathcal{P}\,S^z_\mu(\bq)\,\mathcal{P}^{-1}=-S^z_\mu(\bq),
 \qquad
 \mathcal{P}\,\Phi_h\,\mathcal{P}^{-1}=+\Phi_h .
 \label{eq:parity}
\end{equation}
Since the ground state is a parity eigenstate, $S^z(\bq)|0\rangle$ and
$\Phi(\bq)|0\rangle$ lie in opposite parity sectors and can reach no level in
common.

Figure~\ref{fig:parity} confirms this directly.  Alongside $S^{zz}$ we evaluate
the magnetic channel $|\langle n|\Phi(\bq)|0\rangle|^2$ with
$\Phi(\bq)=\sum_h e^{-i\bq\cdot\br_h}\Phi_h$ summed over the $32$ contractible
plaquettes.  Every level is bright in exactly one channel, with the complement
between $10^{-25}$ and $10^{-35}$; at $\Jpm/\Jzz=-0.20$ the lowest $X$ level
carries magnetic weight $1.645$ against longitudinal weight
$1.1\times10^{-27}$, and at $+0.046$ the figures are $2.891$ and
$1.3\times10^{-26}$.  The excitation the longitudinal probe misses is therefore
an ordinary gauge excitation seen plainly in the magnetic channel, not a
missing or anomalous state.

Two consequences follow for how Fig.~\ref{fig:dssf} should be read.  First, the
$X$ block of the longitudinal response is not the lowest excitation at $X$ and
must not be quoted as one.  Second, because no level is reached by both probes,
no single level on this cluster satisfies the criterion of carrying both an
electric and a magnetic component, so we make no photon assignment here.  We
note that only the $32$ contractible hexagons enter $\Phi$: three quarters of
the six-cycles on this cluster wrap the periodic box, and their ring exchange
transports polarization between transport sectors, which is precisely the
winding artifact the mask removes.  Restricted to the contractible set the ring
exchange conserves the sector exactly, which is an independent confirmation of
the construction of Sec.~\ref{sec:method}.

Finally, none of these energies is converged in system size.  Between cubic-16
and FCC-32 the lowest excitation at $X$ moves by a factor $1.5$--$2.2$ and at
$\Gamma$ by $3.5$--$6.1$, and the ordering of the $\Gamma$ and $X$ levels
inverts.  The selection rule (\ref{eq:parity}) is exact and size-independent;
the energies it acts on are not, and we quote them only as the spectrum of the
finite cluster.
